\section{Теорема Хана-Банаха.}
\begin{theorem}[Хан, Банах, частный случай]
    Пусть $X$ -- линейное нормированное пространство, $L_0 \subset X$ -- его подпространство и $f_0\colon L_0 \to \R$ -- линейный непрерывный на $L_0$ функционал, то есть
    \[
        \|f_0\| = \sup_{x \in L_0, \ \|x\| = 1} |f_{0}x| < +\infty.
    \]
    Тогда существует $f\colon X \to \R$ такой, что $\|f\| = \|f_0\|$ и $f|_{L_0} = f_0$.
\end{theorem}
\begin{proof}
    Докажем более общее утверждение.
    Пусть $p(x)$ -- положительно однородный субаддитивный функционал на $X$, то есть
    \[
        \forall \lambda \geq 0 \hookrightarrow p(\lambda x) = \lambda p(x), \quad p(x + y) \leq p(x) + p(y).
    \]
    И для функционала $f_0$ на $L_0$ справедливо, что
    \[
        \forall x \in L_0 \hookrightarrow f_0(x) \leq p(x).
    \]
    Покажем, что тогда существует продолжение на $X$ с сохранением мажоранты $p$, то есть такой функционал $f\colon X \to \R$, что
    \[
        f|_{L_0} = f_0, \quad \forall x \in X \hookrightarrow f(x) \leq p(x).
    \]
    Тогда утверждение теоремы получается применением этого обобщения к $p(x) = \|f_0\|\|x\|$.

    Рассмотрим $\Phi$ -- семейство линейных продолжений с сохранением мажоранты:
    \[
        \Phi = \left\{(L, g) \mid L_0 \subset L \subset X, \  g\colon L \to \R, \ g|_{L_0} = f_0, \ \forall x \in L \hookrightarrow g(x) \leq p(x)\right\}.
    \]
    Заметим, что $\Phi \neq \emptyset$, поскольку $(L_0, f_0) \in \Phi$.
    Упорядочим $\Phi$ по вложению и совпадению продолжений, то есть
    \[
        (L_1, f_1) \leq (L_2, f_2) \Longleftrightarrow L_1 \subset L_2 \wedge f_2|_{L_1} = f_1.
    \]

    Тогда, согласно теореме Хаусдорфа, в $\Phi$ имеется максимальная цепь $C$.
    Построим мажоранту этой цепи.
    Пусть
    \[
        L^* = \bigcup_{(L, g) \in C} L.
    \]
    $L^*$ -- линейное подпространство в $X$, в силу того, что всякие два элемента $L^*$ принадлежат некоторым $L_\alpha, L_\beta$ из $C$ и, поскольку они должны быть сравнимы, оба элемента лежат либо в $L_\alpha$, либо в $L_\beta$,
    а значит и их линейная комбинация лежит там.
    Определим $g^*\colon L^* \to \R$ как $g^*|_{L} = g$ для всякого $(L, g)$ из $C$.
    Он определён корректно, поскольку всякие два элемента цепи -- сравнимы и функционалы совпадают на пересечении доменов.
    Нетрудно заметить, что для этого функционала выполняется и сохранение мажоранты $p$: всяких $x$ лежит в некотором $L_\alpha$, а для него условие сохранение мажоранты выполнено.
    Таким образом $(L^*, g^*) \in \Phi$ и, в силу максимальности цепи, лежит и в $C$.

    \begin{notet}
        Можно провести рассуждение и с использованием леммы Цорна.

        Пусть $C \subset \Phi$ -- некоторая цепь, и она имеет вид
        \[
            C = \left\{ (L_\alpha, g_\alpha) \right\}_{\alpha \in I},
        \]
        где $I$ -- некоторое индексное множество.
        Рассмотрим
        \[
            L^* = \bigcup_{\alpha \in I} L_{\alpha}.
        \]
        Оно, очевидно, также является линейным подпространством $X$.
        Определим теперь функционал $g^*\colon L^* \to \R$ как
        \[
            g^*(x) = g_{\alpha}(x), \quad x \in L_{\alpha}.
        \]
        Это определение корректно, поскольку любые два элемента цепи сравнимы и функционалы, им соответствующие, совпадают на пересечении доменов.
        Следовательно, $(L^*, g^*) \in \Phi$ и, более того, $(L_\alpha, g_\alpha) \leq (L^*, g^*)$ при всяком $\alpha$.
        Поэтому мы имеем право применять лемму Цорна и в любой цепи находить максимальный элемент.

    \end{notet}

    Покажем, зачем мы всё это делали.
    Пусть $(L, g) \in \Phi$ -- некоторое расширение $f_0$.
    Предположим, что $L \neq X$ -- поскольку иначе мы уже победили.
    Тогда существует $z \notin L$.
    Рассмотрим $\hat{L}$ -- подпространство, натянутое на $L$ и $z$.
    Для определения расширения $g^*$ на $\hat{L}$ нам достаточно задать его на $z$ -- в силу линейности:
    \[
        g^*(a z + y) = a g^*(z) + g^*(y) = a g^*(z) + g(y).
    \]
    Пусть $\alpha, \beta > 0$ и $y_1, y_2 \in L$.
    Тогда
    \begin{multline*}
        \beta g(y_1) + \alpha g(y_2) = g(\beta y_1 + \alpha y_2) = (\alpha + \beta) g\left( \frac{\beta}{\alpha + \beta}y_1 + \frac{\alpha}{\alpha + \beta} y_2\right) \leq (\alpha + \beta) p\left( \frac{\beta}{\alpha + \beta}y_1 + \frac{\alpha}{\alpha + \beta} y_2\right) = \\
        = (\alpha + \beta)p\left(\frac{\beta}{\alpha + \beta}(y_1 - \alpha z) + \frac{\alpha}{\alpha + \beta} (y_2 + \beta z)\right) \leq \beta p(y_1 - \alpha z) + \alpha p(y_2 + \beta z).
    \end{multline*}
    Таким образом, для всяких $\alpha, \beta > 0$ и $y_1, y_2 \in L\colon$
    \[
        \frac{1}{\alpha}\left(-p(y_1 - \alpha z) + g(y_1)\right) \leq \frac{1}{\beta} \left( p(y_2 + \beta z) - g(y_2) \right).
    \]
    Значит, мы можем найти такое $a \in \R$, что
    \[
        \sup_{y \in L, \ \alpha > 0} \left(\frac{1}{\alpha}(-p(y - \alpha z) + g(y))\right) \leq a \leq \inf_{y \in L, \ \alpha > 0} \left(\frac{1}{\alpha}(p(y + \alpha z) - g(y))\right).
    \]
    Но если мы примем $g^*(z) = a$, то подобное продолжение будет удовлетворять вот такому неравенству:
    \[
        - p(y - \alpha z) + g(y) \leq  \alpha g^*(z) \leq p(y + \alpha z) - g(y).
    \]
    Но это в точности означает сохранение мажоранты для $g^*$ на $\hat{L}$.
    Следовательно, теперь мы умеем продолжать функционал на пространство размерности выше на единицу.

    А значит, для максимальной цепи $C$ -- для которой мы построили мажоранту -- мы построили продолжение на большее подпространство.
    Оно будет сравнимо со всяким элементом цепи и, значит, цепь не будет максимальной -- противоречие.
    Значит $L = X$ и мы нашли искомый функционал.

    \begin{notet}
        И закончить рассуждение с использованием леммы Цорна.

        Согласно лемме Цорна, у нас существует максимальный элемент в $\Phi$.
        Предположим, что его область определения меньше $X$.
        Но, тогда, мы можем его продолжить на пространство размерности выше -- а следовательно он не будет максимальным.
        Поэтому его область определения -- $X$, и он является искомым функционалом.
    \end{notet}

\end{proof}

% Записать доказательство Хана-Банаха конкретно от Конста

\begin{corollary}
    Пусть $X$ -- нетривиальное линейное нормированное пространство.
    Тогда существует нетривиальный непрерывный линейный функционал.
\end{corollary}
\begin{proof}
    Пусть $x_0 \in X \setminus \{0\}$.
    Тогда, если мы определим $f(x_0) = t \in \R$, то такой функционал будет непрерывен на $\{\alpha x_0 \mid \alpha \in \R\}$
    \[
        \frac{|f(\alpha x_0)|}{\|\alpha x_0\|} = \frac{|\alpha||f(x_0)|}{|\alpha|\|x_0\|} = \frac{|t|}{\|x_0\|} < +\infty.
    \]
    По теореме Хана-Банаха мы можем продолжить его на всё пространство.
\end{proof}

\begin{example}[Пример ненормируемого линейного метрического пространства (линейное пространство с инвариантной относительно сдвигов метрикой)]
    В частности, покажем что на этом пространстве нет нетривиальных непрерывных линейных функционалов.

    Рассмотрим пространство $L^p([0, 1])$, где $p \in (0, 1)$ с метрикой
    \[
        \rho(f, g) = \int_{[0, 1]} (f - g)^p d\mu.
    \]

    Рассмотрим линейный непрерывный функционал $f\colon L^p([0, 1]) \to \R$.
    Пусть $U(0) = f^{-1}((-1, 1))$ (так как $f(0) = 0$ и $0 \in (-1, 1)$, то полученная окрестность будет окрестностью нуля).
    Заметим, что она будет выпуклой: если $x, y \in U(0)$, то и $\forall t \in [0, 1] \hookrightarrow f(tx + (1 - t)y) = t f(x) + (1 - t) f(y) \in (-1, 1)$ -- так как интервал выпуклый -- и, следовательно и $tx + (1 - t)y \in U(0)$.

    Покажем, что единственным выпуклым множеством, содержащим 0 является всё пространство $L^p$.
    Пусть $x \in L^p([0, 1])$.
    Заметим, что $\forall N \in \N$ существует такое разбиение $\{t_n\}_{n = 0}^{N}$ отрезка $[0, 1]$, что
    \[
        \int_{t_n}^{t_{n + 1}} |x|^{p} d\mu = \frac{1}{N} \int_{[0, 1]} |x|^p d\mu,
    \]
    в силу непрерывности интеграла Лебега.
    Рассмотрим теперь $x_k = x \cdot \Ind_{[t_{k - 1}, t_k]} \cdot N$.
    Исходная $x$ является выпуклой комбинацией $x_k\colon$
    \[
        x = \frac{1}{N}\sum_{k} x_k.
    \]
    Рассмотрим расстояния.
    Поскольку
    \[
        \rho(x_k, 0) = \int_{[t_{k - 1}, t_k]} N^p|x|^p d\mu = N^{p - 1} \rho(x, 0).
    \]
    Поскольку $p - 1 < 0$, то $N^{1 - p} < 1$ и, следовательно, расстояние до $\rho(x, 0)$ строго больше чем $\rho(x_k, 0)$ и стремится к 0 при $N \rightarrow \infty$.

    Заметим, что $\exists r > 0$ такой, что $O_r(0) \subset U(0)$.
    Мы можем выбрать $N$ такой, чтобы $\rho(x_k, 0) < r$ и, следовательно, они будут лежать в шаре $O_r(0) \subset U(0)$.
    Но мы предполагаем, что $U$ -- выпукла.
    Следовательно, и её выпуклая комбинация $x$ -- также лежит в $U$.
    Точка $x$ была выбрана произвольно, а значит $U$ совпадает со всем пространством.

    Тогда $f^{-1}((-1, 1)) = L^p([0, 1])$ и, следовательно, для всякого $x \in L^p$ и $\lambda > 0\colon$
    \[
        |f(\lambda x)| < 1 \Rightarrow |f(x)| < \frac{1}{\lambda} \rightarrow 0, \lambda \rightarrow \infty.
    \]
    Поэтому $f \equiv 0$.
\end{example}
\section{Сильная операторная топология в $\LL(X, Y)$.}
\begin{definition}
    Если $X$, $Y$ -- нормированные пространства, то сильной операторной топологией на $\LL(X, Y)$ называется топология, индуцированная топологией Тихонова из $Y^{X}$.

    То есть если мы рассмотрим канонические проекции $\pi_x\colon$
    \[
        \pi_x(A) = A(x), \quad \forall x \in X, \ \forall A \in \LL(X, Y),
    \]
    то это будет слабейшая топология на $\LL(X, Y)$ такая, что все эти проекции -- непрерывны.
\end{definition}
\begin{example}[Поточечная сходимость к разрывному оператору]
    Рассмотрим $X = (l^1, \|\cdot\|_\infty)$ и $Y = \R$.
    Определим $f_n$ как
    \[
        f_n(x) = \sum_{k = 1}^{n} x(k).
    \]
    Каждый такой оператор -- ограничен и $\|f_n\| = n$.

    У нас есть поточечная сходимость к оператору $f$, который действует как
    \[
        f(x) = \sum_{k = 1}^{\infty} x(k).
    \]
    Он действительно корректно определён и сходимость действительно есть, так как
    \[
        |f_n(x) - f(x)| = \left|\sum_{k = n + 1}^{\infty} x(k)\right| \leq \sum_{k = n + 1}^{\infty} |x(k)| \rightarrow 0, n \rightarrow \infty.
    \]
    Но, при этом, $\|f\| = \infty$ -- можно взять последовательность из точек, у которых первые $n$ элементов -- 1, остальные нули.
    Поэтому $f \notin \LL(X, \R)$ и последовательность $f_n$ в $\LL(X, \R)$ не имеет предела.
\end{example}
\begin{question}
    Таким образом возникает вопрос -- а что надо потребовать, чтобы если $\{A_n\} \subset \LL(X, Y) \subset Y^X$ и $A_n \rightarrow T \in Y^X$ поточечно, то $T \in \LL(X, Y)$?
\end{question}
\begin{note}
    Рассмотрим $x \in B_1(0) \subset X$.
    Поскольку
    \[
        \|Tx\| = \lim \|A_n x\| = \liminf \|A_n x\| \leq \liminf \|A_n\|.
    \]
    Следовательно, если $\{\|A_n\|\}$ -- ограничено и $\forall n \in \N \hookrightarrow \|A_n\| \leq M$, то $\|T\| \leq \liminf \|A_n\| \leq M$.

    Посмотрим на условие ограниченности множества норм операторов.
    По сути, оно эквивалентно тому, что при некотором $M$
    \[
        \forall n \in \N \hookrightarrow A_n B_1(0) \subset B_M(0) \subset Y.
    \]
    То есть, переходя к прообразам, мы получаем
    \[
        \forall n \in \N \hookrightarrow B_\frac{1}{M}(0) \subset A_n^{-1} B_1(0).
    \]
    Или, эквивалентно,
    \[
        0 \in \Int \bigcap_{n = 1}^{\infty} A_n^{-1}(B_1(0)).
    \]

    В нормированном пространстве $B_1(0)$ -- выпуклое симметричное множество.
    Значит, и $A_n^{-1}(B_1(0))$ -- также является выпуклым симметричным множеством.
    Поэтому и $\bigcap_{n} A_n^{-1}(B_1(0))$ -- всё ещё является выпуклым симметричным множеством.
    Поэтому, если внутренность пересечения непуста, то есть некоторый $x \in \Int \bigcap_{n = 1}^{\infty} A_n^{-1}(B_1(0))$, то в ней лежит и $-x$, и $\frac{1}{2}x + \frac{1}{2}(-x) = 0$.
    А значит условие ограниченности множества норм операторов эквивалентно тому, что
    \[
        \Int \bigcap_{n = 1}^{\infty} A_n^{-1}(B_1(0)) \neq \emptyset.
    \]
\end{note}
\subsection{Теорема Бэра.}
\begin{theorem}[Бэр]
    Пусть $(Z, \rho)$ -- полное метрическое пространство и
    \[
        Z = \bigcup_{n = 1}^{\infty} S_n.
    \]
    Тогда существует $n_0 \in \N$ такое, что $\Int [S_{n_0}] \neq \emptyset$.
\end{theorem}
\begin{proof}
    Предположим противное -- то есть $\forall n \in \N \hookrightarrow \Int [S_n] = \emptyset$.
    Значит, для $z_0 \in Z$ верно, что $O_1(z_0) \not\subset [S_1]$.
    Мы можем выбрать точку $z_1 \in O_1(z_0) \setminus [S_1]$ такую, что она входит в $O_1(z_0) \setminus [S_1]$ вместе с некоторым $O_{r_1}(z_1)$, причём $r_1 \leq \frac{1}{2}$.
    Для $O_{r_1}(z_1)$ верно, что $O_{r_1}(z_1) \cap [S_1] = \emptyset$.

    Мы можем аналогичным образом продолжить построение и получить последовательность вложенных шаров с центрами $\{z_k\}$ и их радиусов $\{r_k\}$ такую, что $r_k \leq 2^{-k}$ и что $O_{r_k}(z_k) \cap [S_j], \ j = \overline{1, k} = \emptyset$, и при этом каждый раз строя следующий шар как вложенный в $r_n/4$ -- чтобы было выполнено включение $B_{r_n/2}(z_n) \subset O_{r_n}(z_n)$.
    В силу принципа вложенных шаров $\exists! z^* \in \bigcap_{n = 1}^{\infty} B_{r_n/2}(z_n)$.
    Но, тогда, $z^*$ лежит в некотором $S_n$.
    В то же время она находится и в $B_{r_n/2}(z_n)$ -- противоречие, так как по построению $B_{r_n}(z_n) \cap [S_n] = \emptyset$.
\end{proof}
\begin{example}
    Рассмотрим $C[0, 1] \subset [0, 1]^{[0, 1]}$ с топологией Тихонова.
    Покажем, что $\Ind_{\Q} \notin [C[0, 1]]_{seq}$, но $\Ind_{\Q} \in \left[[C[0, 1]]_{seq}\right]_{seq}$

    Пусть существует $\{x_n\} \subset C[0, 1]$ такая, что $x_n \rightarrow \Ind_{\Q}$.
    Рассмотрим $S_n$, определяемые как
    \[
        S_n = \left\{t \in [0, 1] \mid \forall m \geq n \hookrightarrow |x_m|(t) \leq \frac{1}{2}\right\}.
    \]
    Пусть $\Q \cap [0, 1] = \{r_n\}_{n = 1}^{\infty}$ -- некоторая нумерация.
    В силу поточечной сходимости к $\Ind_{\Q}$
    \[
        [0, 1] \setminus \Q = \bigcup_{n = 1}^{\infty} S_n.
    \]
    Тогда
    \[
        [0, 1] = \bigcup_{n = 1}^{\infty} S_n \cup \bigcup_{n = 1}^{\infty} \{r_n\}.
    \]
    Поскольку $[0, 1]$ -- полно, то по теорему Бэра замыкание одного из $S_n$ содержит внутреннюю точку (очевидно, что замыкание $\{r_n\}$ внутренних точек не содержит).
    В тоже время, $S_n$ -- замкнуто, как пересечение прообразов замкнутого множества относительно непрерывного отображения.
    Но всякое $S_n$ состоит из иррациональных точек и, следовательно, не может содержать интервал -- противоречие.

    Очевидно, что функция Дирихле лежит в втором секвенциальном замыкании: мы можем в первом приблизить всякий индикатор рациональной точки, а во втором -- просто просуммировать по этим рациональным точкам.
\end{example}
\subsection{Теорема Банаха-Штейнгауза.}
\begin{theorem}[Банах, Штейнгауз, принцип равномерной ограниченности]
    Пусть $X$ -- банахово пространство, $Y$ -- нормированное.
    Тогда если последовательность $\{A_n\} \subset \LL(X, Y)$ ограничена поточечно, то есть
    \[
        \forall x \in X \ \exists C(x) > 0\ \forall n\ \|A_n x\| \leq C(x),
    \]
    то последовательность норм $\{\|A_n\|\}$ -- ограничена.
\end{theorem}
\begin{note}
    Банаховость $X$ -- существенна: примером может служить $(l^1, \|\cdot\|_\infty)$ и та же самая последовательность функционалов, что и в предыдущей секции.
\end{note}
\begin{proof}
    Рассмотрим $S = \bigcap_{n = 1}^{\infty} A_n^{-1}(B_1(0))$

    Для всякого $x \in X$ существует $N(x) \in \N$ такое, что
    \[
        \{A_n(x)\}_{n} \subset B_{N(x)}(0).
    \]
    То есть
    \[
        \forall n \in \N \hookrightarrow A_n\left(\frac{x}{N(x)}\right) \subset B_1(0).
    \]
    И
    \[
        \frac{x}{N(x)} \in \bigcap_{n = 1}^{\infty} A_n^{-1}(B_1(0)) = S.
    \]
    Значит мы можем исчерпать $X$ множествами вида $NS$ для $N \in \N$, то есть
    \[
        X = \bigcup_{N = 1}^{\infty} NS.
    \]
    По теореме Бэра при некотором $N_0$ внутренность $N_{0}S$ -- непуста ($N_0 S$ -- замкнуто, его замыкание совпадает с ним самим).
    Тогда и внутренность $S$ -- непуста, из чего и следует утверждение теоремы.
\end{proof}
\begin{corollary}
    Если $X$, $Y$ -- банаховы пространства, то $\LL(X, Y)$ -- полно относительно поточечной сходимости, то есть если
    \[
        \{A_n\} \subset \LL(X, Y)\colon \forall x \in X \hookrightarrow \{A_{n}x\} \text{ -- фундаментальна в } Y, \exists T \in \LL(X, Y)\colon A_n \rightarrow T \text{ поточечно}.
    \]
\end{corollary}
\begin{proof}
    В силу полноты $Y$ и поточечной фундаментальности у нас определён линейный оператор $T$ как
    \[
        T(x) = \lim A_n(x),
    \]
    при всяком $x \in X$.

    С другой стороны, поскольку $X$ -- банахово и из поточечной фундаментальности следует и поточечная ограниченность, то у нас и последовательность норм $A_n$ также ограничена:
    \[
        \exists M > 0 \colon \forall n \in \N \hookrightarrow \|A_n\| \leq M.
    \]

    Тогда, поскольку норма $T$ может быть оценена как
    \[
        \|T\| \leq \liminf \|A_n\| \leq M < \infty,
    \]
    то $T \in \LL(X, Y)$ и теорема доказана.
\end{proof}
\begin{note}
    Полноту $Y$ требовать обязательно.

    Тем не менее, от $X$ нам необходимо лишь потребовать, чтобы оно было множеством второй категории Бэра.
    Существуют неполные нормированные пространства второй категории.

    В частности, в любом бесконечномерном банаховом пространстве $X$ существует собственное всюду плотное подпространство $L$, имеющее вторую категорию Бэра.
\end{note}
\section{Теоремы об обратном операторе и открытом отображении.}
\begin{note}
    Пусть $X$, $Y$ -- нормированные пространства, $A \in \LL(X, Y)$ и $\ker A = 0$ -- то есть $A$ является биекцией на свой образ.
    Тогда у нас существует обратный оператор, определённый на образе:
    \[
        A^{-1}\colon \im A \to X.
    \]
    Он является линейным, поскольку если $y_1, y_2 \in \im A$ и $\alpha_1, \alpha_2 \in \R$, то
    \[
        A(\alpha_1 x_1 + \alpha_2 x_2) = \alpha_1 y_1 + \alpha_2 y_2,
    \]
    а значит
    \[
        \alpha A^{-1}y_1 + \alpha_2 A^{-1}y_2 = A^{-1}(\alpha_1 y_1 + \alpha_2 y_2).
    \]

\end{note}
\begin{note}
    Заметим, что из того, что существует обратный оператор к ограниченному вовсе не следует, что он будет непрерывен.
    \begin{example}
        Рассмотрим $(Ax)(t) = \int_0^t x(\tau)d\tau$, где $A\colon C[0, 1] \to C[0, 1]$.

        Образ $A$ -- это $C^1[0, 1] \cap \{x \mid x(0) = 0\}$ и обратный оператор -- оператор дифференцирования.

        Тем не менее, он не является ограниченным -- $\{\sin nt\}$ как пример.
    \end{example}
\end{note}
\begin{note}
    Заметим, что ограниченность обратного оператора можно переформулировать как ограниченность нормы оператора снизу, то есть
    \[
        \|A^{-1}\| < \infty \Rightarrow x = A^{-1}Ax \Rightarrow \|x\| \leq \|A^{-1}\| \|Ax\|.
    \]
    То есть
    \[
        \|Ax\| \geq C\|x\|, \quad  C > 0.
    \]

    Верно и обратное -- если оператор ограничен снизу, то обратный будет непрерывен:
    \[
        \|A^{-1}\| = \sup_{y \neq 0}\frac{\|A^{-1}y\|}{\|y\|} = \sup_{x \neq 0} \frac{\|x\|}{\|Ax\|} \leq \sup_{x \neq 0} \frac{\|x\|}{C \|x\|} = \frac{1}{C} < +\infty.
    \]
\end{note}
\begin{note}
    Заметим, что из этого условия наличия непрерывного обратного можно извлечь полезное следствие.
    В частности, если $X$ -- полно, $A \in \LL(X, Y)$ и $\exists A^{-1} \in \LL(\im A, X)$, то $\im A$ -- полно.
\end{note}
\begin{definition}
    Отображение $f\colon (X, \tau_X) \to (Y, \tau_Y)$ называется открытым, если образ открытого $U \in \tau_X$ является открытым $f(U) \in \tau_Y$.
\end{definition}
\begin{proposition}[Топологический критерий непрерывной обратимости.]
    Пусть $A \in \LL(X, Y)$ и имеет нулевое ядро.

    Тогда $A^{-1} \in \LL(\im A, X) \Longleftrightarrow A$ -- открытое отображение.
\end{proposition}
\begin{proof}
    Пусть $G \subset X$ -- открытое.
    Тогда $(A^{-1})^{-1}(G) = A(G)$ и $A(G) \subset \im A$ -- должно быть открыто для всякого открытого $G$.

    Очевидно верно и обратное.
\end{proof}
\begin{note}
    Заметим, что мы можем переформулировать определение открытого отображения в случае линейного оператора как
    \[
        \exists r > 0\colon O_r(0) \cap \im A \subset A(O_1(0)).
    \]

    Очевидно, что если $A$ -- открытый оператор, то $A(O_1(0))$ -- открытый и $0$ лежит в образе вместе с некоторым шаром.

    Пусть $\exists r > 0 \colon O_r(0) \cap \im A \subset A(O_1(0))$.
    Покажем, что образом открытого является открытое.
    Пусть $G \subset X$ -- открытое.
    Рассмотрим некоторую точку $x \in G$, она входит в $G$ вместе с некоторым шаром $O_R(x)$.
    Заметим, что
    \[
        O_R(x) = x + RO_1(0).
    \]
    Значит, существует $r > 0$ такое, что
    \[
        O_r(0) \cap \im A \subset A(O_1(0)) \Rightarrow O_{Rr}(x) \cap \im A \subset A(O_R(x)).
    \]
    И точка $Ax$ также является внутренней -- следовательно, оператор открыт.
\end{note}
\begin{theorem}[Банах, об открытом отображении]
    Пусть $X$, $Y$ -- банаховы пространства, $A \in \LL(X, Y)$ и $\im A = Y$.

    Тогда $A$ -- открытое отображение.
\end{theorem}
\begin{proof}
    Рассмотрим
    \[
        \bigcup_{n = 1}^{\infty} nA(O_1(0)) = \bigcup_{n = 1}^{\infty} A(O_n(0)) = Y.
    \]
    Согласно теореме Бэра существует $n_0 \in \N\colon$
    \[
        \Int [A(O_{n_0}(0))] \neq \emptyset,
    \]
    то есть и
    \[
        \Int [A(O_1(0))] \neq \emptyset.
    \]
    Заметим, что $[AO_1(0)]$ -- выпукло и симметрично относительно нуля.
    Поэтому, если есть внутренняя точка $x$, то и $-x$ -- также внутренняя точка, а следовательно и 0 будет внутренней точкой замыкания, то есть $\exists r_0\colon$
    \[
        O_{r_0}(0) \subset [AO_1(0)].
    \]

    До этого мы никак не использовали полноту $X$.
    Для всякого $\epsilon > 0$ верно следующее включение:
    \[
        O_{\epsilon r_0}(0) \subset [AO_{\epsilon}(0)].
    \]

    Для любой $y \in [AO_1(0)]$ множество
    \[
        y - \left[AO_{\frac{1}{2}}(0)\right]
    \]
    содержит некоторую окрестность $y$ -- поскольку у нас во всяком замыкании образа шара есть некоторая окрестность нуля.

    Более того, её пересечении с $AO_1(0)$ -- непусто, так как $y$ лежит в замыкании.
    Значит существует $x_1 \in O_1(0)$ и $y_1 \in \left[AO_{\frac{1}{2}}(0)\right]$ такие, что
    \[
        y - y_1 = Ax_1.
    \]

    Аналогичным образом мы делаем индукционный переход -- то есть от $y_1$ мы отнимаем шар радиуса $\frac{1}{4}$, и так далее, и тому подобное.
    Таким образом мы построим две последовательности -- $\{y_n\}$ такую, что $y_n \in [AO_{2^{-n}}(0)]$ и $\{x_n\}$ такую, что $y_n - y_{n + 1} = Ax_{n + 1}$ и $x_n \in O_{2^{-n + 1}}(0)$.

    Тогда
    \[
        y - y_{n + 1} = y - y_1 + y_1 - \ldots - y_{n - 1} = Ax_1 + \ldots Ax_{n + 1}.
    \]
    Если $S_n = \sum_{k = 1}^{n}x_n$, то $y - y_{n + 1} = AS_{n + 1}$.

    Заметим, что
    \[
        \sum_{n = 1}^{\infty} \|x_n\| \leq \sum_{n = 1}^{\infty} 2^{-n + 1} = 2.
    \]
    Значит и $S_n \rightarrow x$ -- в полном пространстве абсолютно сходящийся ряд сходится и $x \in O(3)$.

    В силу непрерывности оператора $A$
    \[
        \|y_n\| \leq 2^{-n}\|A\| \rightarrow 0.
    \]
    А значит
    \[
        y = Ax \subset AO_3(0).
    \]
    Таким образом
    \[
        AO_1(0) \supset O_{\frac{r_0}{3}}(0),
    \]
    из чего и следует открытость оператора.
\end{proof}
\begin{theorem}[Банах, об обратном операторе]
    Пусть $X$, $Y$ -- банаховы пространства, $A \in \LL(X, Y)$ и является биекцией.

    Тогда $A^{-1} \in \LL(Y, X)$.
\end{theorem}
\begin{proof}
    Очевидно.
\end{proof}